\section{Definici\'on y tipos de requisitos}
El handbook define \textbf{requisito} desde tres puntos de vista complementarios: como una necesidad percibida por un stakeholder, como una capacidad o propiedad que el sistema debe tener y como la representaci\'on documentada de esa necesidad o capacidad \cite[p.~10]{cpre}.

Con esa base, el documento insiste en que los requisitos existen incluso cuando no se han documentado, y que forman la base de cualquier desarrollo o evoluci\'on de sistemas. Tambi\'en diferencia tres clases principales: requisitos funcionales, requisitos de calidad y restricciones \cite[p.~10]{cpre}.

El propio handbook ampl\'ia el marco conceptual con la definici\'on de Requirements Engineering como un enfoque sistem\'atico y disciplinado para especificar y gestionar requisitos con el objetivo de entender los deseos y necesidades de los stakeholders y reducir el riesgo de entregar un sistema inadecuado \cite[pp.~10--11]{cpre}. A la vez, define stakeholder como cualquier persona u organizaci\'on que influye en los requisitos o se ve afectada por el sistema \cite[p.~11]{cpre}.